\documentclass[10pt]{article}
% Math Packages
\usepackage{amsmath, mathtools}
\usepackage{amssymb}
\usepackage{amsthm}
\usepackage{amsfonts}
\usepackage{bbm}
\usepackage{breqn}
\usepackage[margin=1in]{geometry}
\usepackage{graphicx}
\usepackage{tikz}
\usetikzlibrary{arrows.meta}
\usetikzlibrary{calc}
\usepackage{forest}
\usepackage{tikz-qtree}
\graphicspath{ {./images/} }
\usepackage{hyperref}
\usepackage[capitalize]{cleveref}
\usepackage[shortlabels]{enumitem}
\usetikzlibrary{arrows,matrix,positioning}
\usepackage{multicol}


\pagenumbering{gobble} % Removes page numbers

% for the pipe symbol
\usepackage[T1]{fontenc}

% Citing theorems by name. (source: https://tex.stackexchange.com/questions/109843/cleveref-and-named-theorems)
\makeatletter
\newcommand{\ncref}[1]{\cref{#1}\mynameref{#1}{\csname r@#1\endcsname}}

\def\mynameref#1#2{%
  \begingroup
  \edef\@mytxt{#2}%
  \edef\@mytst{\expandafter\@thirdoffive\@mytxt}%
  \ifx\@mytst\empty\else
  \space(\nameref{#1})\fi
  \endgroup
}
\makeatother

% Colorful Notes
\usepackage{color} \definecolor{Red}{rgb}{1,0,0} \definecolor{Blue}{rgb}{0,0,1}
\definecolor{Purple}{rgb}{.5,0,.5} \def\red{\color{Red}} \def\blue{\color{Blue}}
\def\gray{\color{gray}} \def\purple{\color{Purple}}
\newcommand{\rnote}[1]{{\red [#1]}} % \rnote{foo} gives '[foo]' in red
\newcommand{\pnote}[1]{{\purple [#1]}} % \pnote{foo} gives '[foo]' in purple
\newcommand{\bnote}[1]{{\blue #1}} % \bnote{foo} gives 'foo' in blue
\newcommand{\gnote}[1]{{\gray #1}} % \gnote{foo} gives 'foo' in gray
\newcommand{\Max}[1]{{\purple [#1]}} % \bnote{foo} then 'foo' is blue


% Claim numbering (the counter restarts after each proof environment)
\newcounter{claimcount}
\setcounter{claimcount}{0}
\newenvironment{claim}{\refstepcounter{claimcount}\par\addvspace{\medskipamount}\noindent\textbf{Claim \arabic{claimcount}:}}{}
\usepackage{etoolbox}
\AtBeginEnvironment{proof}{\setcounter{claimcount}{0}}
\newenvironment{claimproof}{\par\addvspace{\medskipamount}\noindent\textit{Proof of Claim  \arabic{claimcount}.}}{\hfill\ensuremath{\qedsymbol} \tiny{Claim}

  \medskip}
% Add claim support to cleverref
\crefname{claimcount}{Claim}{Claims}


% Math Environments
\newtheorem{theorem}{Theorem}
\newtheorem{assumption}[theorem]{Assumption}
\newtheorem{lemma}[theorem]{Lemma}
\newtheorem{proposition}[theorem]{Proposition}
\newtheorem{corollary}[theorem]{Corollary}
\newtheorem{question}[theorem]{Question}
\theoremstyle{definition}
\newtheorem{definition}[theorem]{Definition}
\newtheorem{remark}[theorem]{Remark}
\newtheorem{example}[theorem]{Example}
\newtheorem{notation}[theorem]{Notation}
\newtheorem{problem}[theorem]{Problem}

% Matrices and Column Vectors. 
\usepackage{stackengine}
\setstackgap{L}{1.0\normalbaselineskip}
\usepackage{tabstackengine}
\setstackEOL{;}% row separator
\setstackTAB{,}% column separator
\setstacktabbedgap{1ex}% inter-column gap
\setstackgap{L}{1.5\normalbaselineskip}% inter-row baselineskip
\let\nmatrix\bracketMatrixstack  %Usage: \nmatrix{1,2,3\4,5,6}
\newcommand\cv[1]{\setstackEOL{,}\bracketMatrixstack{#1}} %usage: \cv{1,2,3}

% Custom Math Coqmmands
\newcommand{\vt}{\vskip 5mm} % vertical space
\newcommand{\fl}{\noindent\textbf} % first line
\newcommand{\Fl}{\vt\noindent\textbf} % first line with space above
\newcommand{\norm}[1]{\left\lVert#1\right\rVert} % norm
\newcommand{\pnorm}[1]{\left\lVert#1\right\rVert_p} % p-norm
\newcommand{\qnorm}[1]{\left\lVert#1\right\rVert_q} % q-norm
\newcommand{\1}[1]{\textbf{1}_{\left[#1\right]}} % indicator function 
\def\limn{\lim_{n\to\infty}} % shortcut for lim as n-> infinity
\def\sumn{\sum_{n=1}^{\infty}} % shortcut for sum from n=1 to infinity
\def\sumkn{\sum_{k=1}^{n}} % shortcut for sum from k=1 to n
\def\sumin{\sum_{i=1}^{n}} % shortcut for sum from i=1 to n
\def\SAs{\sigma\text{-algebras}} % shortcut for $\sigma$-algebras
\def\SA{\sigma\text{-algebra}} % shortcut for $\sigma$-algebra
\def\Ft{\mathcal{F}_t} % time-indexed sigma-algebra (t)
\def\Fs{\mathcal{F}_s} % time-indexed sigma-algebra (s)
\def\F{\mathcal{F}} % sigma-algebra
\def\G{\mathcal{G}} % sigma-algebra
\def\R{\mathbb{R}} % Real numbers
\def\N{\mathbb{N}} % Natural numbers
\def\Z{\mathbb{Z}} % Integers
\def\E{\mathbb{E}} % Expectation
\def\P{\mathbb{P}} % Probability
\def\Q{\mathbb{Q}} % Q probability
\def\dist{\text{dist}} %Text 'dist' for things like 'dist(x,y)'
\newcommand{\indep}{\perp \!\!\! \perp}  %independence symbol
\def\Var{\mathrm{Var}} % Variance
\def\tr{\mathrm{tr}} % trace

% Brackets and Parentheses
\def\[{\left [}
    \def\]{\right ]}
% \def\({\left (}
%   \def\){\right )}



\usepackage{color}
\definecolor{Red}{rgb}{1,0,0}
\definecolor{Blue}{rgb}{0,0,1}
\definecolor{Purple}{rgb}{.5,0,.5}
\def\red{\color{Red}}
\def\blue{\color{Blue}}
\def\gray{\color{gray}}
\def\purple{\color{Purple}}
\definecolor{RoyalBlue}{cmyk}{1, 0.50, 0, 0}
\newcommand{\dempfcolor}[1]{{\color{RoyalBlue}#1}} 
\newcommand{\demph}[1]{\dempfcolor{{\sl #1}}}

% comment exactly one of the following line to show / hide the solutions
% \newcommand{\solution}[1]{{\purple #1}} % uncomment to show the solutions
\newcommand{\solution}[1]{} % uncomment to hide the solutions



\title{Lecture Notes for Math 372: \\Elementary Probability and Statistics}
\date{Last updated: \today}
% \author{mh}

\begin{document}


% \begin{itemize}
  
%   \item The exam will be on Monday, March 24. It will be in-class, so you will
%   have about 50 minutes to do the problems.
%   \item The exam will be closed book with no calculators or notes. 
%   \item The exam will cover material from the first 7 weeks of class (i.e.,
%   lectures 1 through 19). This means that quadric surfaces are not on the exam.
%   \item The practice exams are very similar to the exam. I strongly encourage
%   you to work through the practice exams in a timed environment.
%   \item The exam is meant to be \textit{predictable:} all but one of the
%   problems will be drawn from homework problems (possibly with minor
%   alterations). The other problem is drawn verbatim from the lecture examples.
%   \item I \textit{won't} ask you to do something like Problem 9 from HW 3 or
%   Problem 11 from HW 4. These problems are too difficult for a 50-minute exam.
  
% \end{itemize}



\begin{center}
  \section*{Math 243: Practice Midterm \#1}
\end{center}

\textit{Instructions: you must show your work to recieve credit. You have 50
  minutes. Calculators are not allowed.}

\begin{problem}[Differentiation of curves]
  Consider the curve
  \begin{equation*}
    \left\{ \begin{array}{l@{}l} x(t) = 5t^{2}-6t+4 \\y(t)= t^{2}+6t-1 \end{array}\right.
  \end{equation*}
  \begin{enumerate}[(a)]
    \item Find an equation of the tangent line at $t=0$.
    \item For what values of $t$ is this curve concave down? For what values is it concave up?
  \end{enumerate}
\end{problem}


\begin{problem}[Arc length]
  Compute the arc length of the parametrized curve:
  \begin{equation*}
    \left\{ \begin{array}{l@{}l} x(t) = \frac{4}{3}t^{3}+t-2 \\ y(t) = \frac{1}{27}(9t^{2}+2)^{\frac{3}{2}} \end{array}\right. \qquad 0 \leq t \leq 2
  \end{equation*}
\end{problem}


\begin{problem}[Polar coordinates]
  Consider the polar graph
  \begin{equation*}
    \left\{ \begin{array}{l@{}l} r= 6\sin(\theta) \\ 0 \leq  \theta \leq \pi \end{array}\right.  %
  \end{equation*}
  Convert this equation to Cartesian form and plot it.
\end{problem}

\begin{problem}[Projections]
  Let $\vec{u}= \left\langle 1,-1,1 \right\rangle $ and $\vec{v} = \left\langle -2,-2,-2 \right\rangle $.
  \begin{enumerate}[(a)]
    \item Find the projection of $\vec{u}$ along $\vec{v}$. 
    \item What is the cosine of the angle between $\vec{u}$ and $\vec{v}$?
    \item Find a vector which is orthogonal to $\vec{u}$ and $\vec{v}$.
  \end{enumerate}
\end{problem}

\begin{problem}[Planes]
  Consider the points  $A=(2,5,-1)$, $B=(-2,1,3)$ and $C = (7,1,1)$.
  \begin{enumerate}[(a)]
    \item Find an equation of the line passing through points $A$ and $B$.
    \item Find an equation the plane passing through points $A$, $B$, and $C$.
  \end{enumerate}
\end{problem}


\begin{problem}[Sphere]
  \begin{enumerate}[(a)]
    \item []
    \item State the Cartesian formula for a sphere.
    \item Show that the equation $x^{2}+y^{2}+z^{2} +4x - 6y + 2z +6=0 $ represents a sphere.
    \item What is the center and radius of the sphere?
  \end{enumerate}
\end{problem}

\newpage
\begin{center}
  \section*{Math 243: Practice Midterm \#2}
\end{center}

\textit{Instructions: you must show your work to recieve credit. You have 50
  minutes. Calculators are not allowed.} \setcounter{theorem}{0}

\begin{problem}[Vectors]
  \begin{enumerate}[(a)]
    \item []
    \item Compute the distance between the following two points $(5,8,15)$ and $(1,4,8)$.
    \item Are the vectors $\left\langle 2,4,7 \right\rangle $ and $\left\langle 4,8,14 \right\rangle$
    orthogonal, parallel, or neither? Justify your answer.
    \item Are the vectors $\left\langle 3,-1,3 \right\rangle $ and $\left\langle 5,9,-2 \right\rangle $
    orthogonal, parallel, or neither? Justify your answer.
  \end{enumerate}
\end{problem}


\begin{problem}[Arc length]
  Find the arc length of the parametric curve
  \begin{equation*}
    \left\{ \begin{array}{l@{}l} x(t) = \frac{1}{2}t^{2}-t-1 \\ y(t)=\frac{4}{3}t^{\frac{3}{2}} \end{array}\right.
  \end{equation*}
  on the itnerval $1 \leq t \leq 2$
\end{problem}



\begin{problem}[Polar stuff]
  % [Peanut curve]
  Find the area enclosed by the curve
  \begin{equation*}
    \left\{ \begin{array}{l@{}l} r = 1+ \cos(2\theta) \\ 0 \leq \theta \leq \pi\end{array}\right.
  \end{equation*}
  \begin{center}
    \includegraphics[scale=.4]{images/plot}
  \end{center}
\end{problem}


\begin{problem}[Projections]
  Consider the point $W=(1,0,3)$ and the plane given by the equation $5x-y-z=3$.
  Find the shortest distance from $W$ to the plane. Include a relevant sketch in
  your answer.
\end{problem}


\begin{problem}%[Differentiation of Curves]
  Let $C$ be the curve defined by the parametric equations
  \begin{equation*}
    \left\{ \begin{array}{l@{}l} x(t)=t^{2}-t+2 \\ y(t) = t-\frac{1}{3}t^{3}, \end{array}\right.
  \end{equation*}
  where $-10 \leq t \leq 10$.
  \begin{enumerate}[(a)]
    \item Find the equation of the tangent line to $C$ at $t=2$.
    \item For what values of $t$ is the tangent line horizontal? For what values is it vertical?
    \item Compute $\frac{d^{2}y}{dx^{2}}$.
    \item Is $C$ concave up or concave down when $t=1$? You must justify your answer to recieve credit.
  \end{enumerate}
\end{problem}

\begin{problem}[Planes and cross product]
  The planes given by the equations
  \begin{equation*}
    4x-y+3z = 1 \quad \text{and} \quad  2x-2y+z=1
  \end{equation*}
  intersect to form a line $L$. 
  \begin{enumerate}[(a)]
    \item Find a vector in the direction of the line $L$.
    \item Give an example of a point on the line $L$.
    \item Find an equation of the line $L$ (either a vector equation or a parametric equation is fine).
  \end{enumerate}
\end{problem}


\end{document}